\unitopen{19}
\begin{candos}
\item porównać dwie rzeczy (\de{größer als}, \de{so schnell wie}) i wskazać najlepszą (\de{am besten});
\item znać nieregularne stopniowanie: \de{gut/besser/am besten}, \de{viel/mehr/am meisten}, \de{gern/lieber/am liebsten};
\item opisać pogodę i pory roku;
\item porównać miasta, zespoły i narzędzia prostymi zdaniami.
\end{candos}

\section{Słownictwo}
\vocabtable{19}

\section{Stopniowanie}
\begin{center}\small
\begin{tabular}{@{}llll@{}}
\toprule
\textsf{\textbf{Podstawowy}} & \textsf{\textbf{Wyższy}} & \textsf{\textbf{Najwyższy}} & \textsf{\textbf{Uwaga}}\\
\midrule
\de{schnell} & \de{schneller} & \de{am schnellsten} & regularnie: \de{-er}, \de{am -sten}\\
\de{klein} & \de{kleiner} & \de{am kleinsten} & \\
\de{groß} & \de{größer} & \de{am größten} & przegłos\\
\de{alt} & \de{älter} & \de{am ältesten} & przegłos, \de{-e-}\\
\de{hoch} & \de{höher} & \de{am höchsten} & \de{c} wypada\\
\de{gut} & \de{besser} & \de{am besten} & nieregularny\\
\de{viel} & \de{mehr} & \de{am meisten} & nieregularny\\
\de{gern} & \de{lieber} & \de{am liebsten} & nieregularny\\
\bottomrule
\end{tabular}
\end{center}
Porównanie: \de{Zürich ist größer \textbf{als} Basel.} Równość: \de{Der Server ist \textbf{so} schnell \textbf{wie} der Laptop.}
Przed rzeczownikiem przymiotnik bierze końcówki (jednostka~21): \de{der schnellere Server}. Pogoda zwykle bez podmiotu: \de{es}:
\de{Es regnet. Es ist kalt.}

\begin{falle}
Po polsku: „większy \emph{niż}”, ale po niemiecku zawsze \de{als}, a nie \de{wie} (\de{größer wie} to błąd
potoczny). Przy równości: \de{so \dots{} wie}, nie \de{als}. Drugi błąd: „najlepszy” w roli orzecznika to \de{am besten},
a przed rzeczownikiem \de{der beste Kollege}. Trzeci: \de{viel} nie ma wyższego \de{vieler}: \de{mehr}.
\end{falle}

\begin{szwajcaria}
Szwajcaria to kraj gór i jezior: \de{der Zürichsee}, \de{die Alpen}, \de{der Säntis}. Pogoda w Zurychu jest zmienna;
zimą często jest szaro (\de{Hochnebel}). Temperatury podaje się w stopniach Celsjusza, a o pogodzie rozmawia się
zawsze: \de{Schönes Wetter heute, oder?} to w CH klasyczna zagajka \de{Smalltalk}.
\end{szwajcaria}

\begin{praca}
Porównanie rozwiązań: \de{Die neue Datenbank ist schneller als die alte.} \de{Das ist die beste Lösung.}
\de{Für uns ist Azure wichtiger als AWS.} \de{Wir haben mehr Tests als letztes Jahr.} W code review:
\de{Das ist lesbarer als vorher.}
\end{praca}

\section{Dialogi}
\textbf{Dialog 1 --- pogoda i miasta}
\begin{dialog}
\turn{A}{Wie ist das Wetter in Zürich im Winter?}{Jaka jest pogoda w Zurychu zimą?}
\turn{B}{Es ist kälter als in Spanien und oft neblig. Im Sommer ist es viel wärmer.}{Jest zimniej niż w Hiszpanii i często mgliście. Latem jest dużo cieplej.}
\turn{A}{Welche Jahreszeit magst du am liebsten?}{Którą porę roku lubisz najbardziej?}
\turn{B}{Ich mag den Herbst am liebsten. Das Wetter ist nicht so heiß wie im Sommer.}{Najbardziej lubię jesień. Pogoda nie jest tak gorąca jak latem.}
\end{dialog}

\section{Ćwiczenia}
\begin{cwiczenia}
\exitem{Wyższy stopień: \de{schnell}.}{\de{schneller}}
\exitem{Wyższy stopień: \de{groß}.}{\de{größer}}
\exitem{Wyższy stopień: \de{gut}.}{\de{besser}}
\exitem{Najwyższy stopień: \de{alt} (po \de{am}).}{\de{am ältesten}}
\exitem{Najwyższy stopień: \de{viel}.}{\de{am meisten}}
\exitem{Uzupełnij: \de{Zürich ist größer \dots{} Bern.}}{\de{als}}
\exitem{Uzupełnij: \de{Der Laptop ist so schnell \dots{} der Server.}}{\de{wie}}
\exitem{Uzupełnij: \de{Ich esse \dots{} Suppe als Salat.} (gern, wyższy)}{\de{lieber}}
\exitem{Popraw: \de{Er ist größer wie ich.}}{\de{Er ist größer als ich.}}
\exitem{Opisz pogodę: „Pada deszcz i jest zimno.”}{\de{Es regnet und es ist kalt.}}
\exitem{Przetłumacz: „Ta baza danych jest szybsza niż stara.”}{\de{Diese Datenbank ist schneller als die alte.}}
\exitem{Przetłumacz: „Najbardziej lubię lato.”}{\de{Ich mag den Sommer am liebsten.}}
\end{cwiczenia}

\section{Trening słownictwa (generowany)}
\drill{19}

\weekend{Porównaj dwa miasta (np.\ Zürich i Twoje) w ośmiu zdaniach: pogoda, ceny, wielkość, praca,
z użyciem \de{als}, \de{so \dots{} wie}, \de{am besten}. Nagraj i sprawdź.}
